\vspace{-3mm}
\section{Introduction}
\label{sec:intro}
\vspace{-0.6mm}



Controlling motion dynamics in video generation has received increasing attention~\cite{CameraCtrl,MotionCtrl,HumanVid,AnimateDiff,VMC,STDF,MotionClone,DreamMotion,MotionInversion,MotionMaster,shuai2024survey}, as it enables better customization and is crucial in many applications. For example, in filmmaking, directors meticulously choreograph the movements of both actors and the camera.
Consequently, precise control over object and camera motions in video offers creative flexibility.





Despite recent progress, challenges remain in motion control of text-to-video~(T2V) generation. A key limitation is \textit{\textbf{the lack of high-quality datasets with comprehensive 6D pose annotations}}.
For controlling object movement \cite{MotionCtrl,AnyI2V,VideoComposer,DragNUWA,Tora,Trailblazer,FreeTraj,Motion-Zero,MotionBooth,Image-Conductor,Motion-I2V}, the motion is primarily annotated as the trajectory in image space~\cite{MotionCtrl,DragNUWA}. This annotation, however, lacks 3D nature and intertwines the dynamics of both objects and the camera. For example, a rightward trajectory could represent either a stationary camera with a moving object or a static object with a left-moving camera. Recently, \ang{360}-Motion synthetic dataset \cite{3DTrajMaster} provides 6D poses of objects, but limited with static camera setting and motion diversity.
Besides, existing commonly used datasets~\cite{Direct-A-Video,RealEstate-10K,MVImageNet} for learning camera motion mainly focus on scenes with minimal object dynamics.
Some human-centric synthetic datasets \cite{HumanVid,BEDLAM,SynBody,WHAC} provide ground truth for both human subjects and camera motions within a global coordinate system, yet exhibit limited motion diversity and category variety. 

Another limitation is \textbf{\textit{the absence of methods that can independently or jointly control the 6D poses of both object and camera}}. For example, methods like Motion-Zero \cite{Motion-Zero} animate objects without 3D-aware control (\eg, orientation) \cite{PEEKABOO,MotionBooth,Trailblazer}, while methods like CameraCtrl \cite{CameraCtrl} exclusively focus on camera motion~\cite{Training-free-CC,CamCo,VD3D}.~MotionCtrl \cite{MotionCtrl} trains separate modules for object and camera control in a two-stage process.~However, without access to video data containing complete 6D pose annotations for both elements, it struggles to achieve realistic, synchronized control of objects and cameras within the same scene.

\begin{figure}
    \centering
    \includegraphics[width=0.999\linewidth]{personalized_compressed.pdf}
    \vspace{-7.6mm}
    \caption{Example videos generated by our method \methodname trained on the proposed \datasetname dataset, showing its adaptability  with different personalized T2I models~\cite{StableDiffusion,Toonyou,RealisticVision}.}
    \label{fig:personalized}
    \vspace{-3.6mm}
\end{figure}



To address these limitations, a dataset with comprehensive 6D pose annotations of objects and camera is desired. However, acquiring such data is challenging and typically requires specialized equipment and expertise. In this work, we introduce a \textbf{\underline{Syn}}thetic dataset for \textbf{\underline{F}}ree-Form \textbf{\underline{M}}otion \textbf{\underline{C}}ontrol (\textbf{\datasetname}). Designed with an emphasis on quality and diversity, the dataset includes a rich array of animated object assets and environment assets across various categories. Furthermore, a rule-based generation algorithm is implemented to create trajectories for both objects and the camera, as shown in \cref{fig:teaser}. This algorithm encompasses basic patterns and simulates challenging cinematographic shots as in \cref{fig:curves}. {Compared to recent works~\cite{3DTrajMaster,HumanVid} that can only construct uncontrolled trajectories, our rule-based algorithm supports customized object~\&~camera motions with diverse patterns.} To enhance realism, essential attributes of objects, like living environment, types of speed and size, are annotated using a Multimodal Large Language Model (MLLM)~\cite{InternVL} and manual labeling, facilitating the generation of plausible trajectories. The \datasetname dataset also provides detailed annotations, including 6D pose information of objects and camera, instance segmentation maps, depth maps, and comprehensive descriptions of content and motion, supporting a wide range of research fields.


To further validate the effectiveness of the proposed \datasetname dataset and support {3D-aware control} in T2V generation, we propose a \methodcompletename (\methodname) method. \methodname mainly includes two components: \camctrlcompletename (\camctrlname) and \objctrlcompletename (\objctrlname). Unlike previous methods~\cite{DragNUWA,MotionCtrl}, {Our approach trained on \datasetname disentangles global (camera) and local (object) dynamics and manipulates the 6D poses of camera and objects}. Furthermore, we adopt Domain LoRA~\cite{LoRA} to prevent model from fitting to rendered style in synthetic data. As shown in \cref{fig:personalized} and \cref{fig:lora}, \methodname effectively mitigates the domain gap and adapts to various personalized Text-to-Image (T2I) models, generating high-fidelity results across diverse styles. In addition, \methodname provides flexible user interfaces for motion control. Users can input trajectories for objects and camera by simply drawing 3D curves or by specifying motion types for each (as detailed in \cref{sec:data_pipeline}), which are used by the rule-based algorithm to generate their trajectories accordingly.

\begin{table*}[t]
    \centering
    \renewcommand\arraystretch{1.06}
    \caption{Comparison of the proposed \textbf{\datasetname} with existing datasets. The object/camera motion pattern columns apply only to synthetic datasets. In addition to offering a rich variety of object categories, \datasetname outperforms in motion pattern variety and controllability with comprehensive pose annotations of camera and objects. In our implementation, we only use 26K subset as training data.}
    \vspace{-3.6mm}
    \footnotesize
    \resizebox{1\textwidth}{!}{
    \begin{tabular}{r|ccc>{\centering\arraybackslash}m{0.14\textwidth}>{\centering\arraybackslash}m{0.131\textwidth}>{\centering\arraybackslash}m{0.122\textwidth}>{\centering\arraybackslash}m{0.122\textwidth}}
    \specialrule{.1em}{.05em}{.05em} 
        \rowcolor{mygray!50}Dataset & Clips  & Source &  Category & \makecell[c]{Object Motion\\ Pattern}& \makecell[c]{Camera Motion\\ Pattern}& \makecell[c]{Camera Pose\\Annotation} &  \makecell[c]{Object Motion\\Annotation}  \\ 
        \hline\hline
        RealEstate10K~\cite{RealEstate-10K} & \ \ 65K  & Real & Real Estate &- & -&  Fitting & \xmarkg \\
        \rowcolor{mygray2!36}MVImgNet~\cite{MVImageNet} & 220K  & Real &  \textbf{Common} & - & -&  Fitting & \xmarkg \\

        VideoHD~\cite{DragNUWA} & \ \ 75K & Real & \textbf{Common} & - & - & \xmarkg  & Optical Flow   \\
        
        \rowcolor{mygray2!36}MotionCtrl~\cite{MotionCtrl} & \textbf{240K} &  Real& \textbf{Common} &-  & - & \xmarkg  & Optical Flow \\
        HumanVid-Real~\cite{HumanVid} & \ \ 20K  & Real & Human & - & - & Fitting  & 2D Human Pose \\

        \arrayrulecolor{gray!90}\hline
         BEDLAM~\cite{BEDLAM} & \ \ 10K  & Synthetic & Human & Limited/Uncontrollable & Static & \textbf{Ground Truth}  & 3D Human Pose  \\
        
        \rowcolor{mygray2!36}SynBody~\cite{SynBody} & \ \ 27K  &Synthetic & Human & Limited/Uncontrollable & Static & \textbf{Ground Truth}  & 3D Human Pose \\
        
        
        HumanVid-Syn~\cite{HumanVid} & 100K  & Synthetic & Human & Limited/Uncontrollable & \textbf{Diverse}/{Uncontrollable} & \textbf{Ground Truth}  & 3D Human Pose \\

        \rowcolor{mygray2!36}\ang{360}-Motion~\cite{3DTrajMaster} &\ \ 54K  & Synthetic & Animal \& Human & Limited/Uncontrollable & Static & \xmarkg  & \textbf{Object Pose} \\
        
        \arrayrulecolor{black}\hline    
        \rowcolor{green!8}\textbf{\datasetname} (\textbf{ours}) & \ \ 62K  & Synthetic & \textbf{Common} & \textbf{Diverse/Controllable} & \textbf{Diverse/Controllable} & \textbf{Ground Truth}  & \textbf{Object Pose} \\        
     \specialrule{.1em}{.05em}{.05em} 
    \end{tabular}
    }
    \label{tab:datasets}
    \vspace{-5.6mm}
\end{table*}


In summary, our main contributions are as follows:
\begin{itemize}
    \item To the best of our knowledge, the \datasetname dataset is the first to provide 6D pose annotations for both camera and objects. Its diverse scenes and complex motion patterns provide models with valuable resources for learning the dynamics of multiple objects and the camera. 

    \item {The \methodname method can manipulate the 6D poses of the camera and objects independently or simultaneously, achieving high-quality results across diverse scenes.}

    \item Extensive experiments demonstrate that \methodname, trained on \datasetname dataset, generates videos of superior quality compared to state-of-the-art methods.

\end{itemize}
